% !TeX root = ../../src/main.tex
% ============================================================================
% ERRATA RESUMO -- o escopo do resultado de categoria no Resumo em portugues.
% Achada 2026-08-21, na preparacao da apresentacao de defesa.
% Destino: o Resumo (pt) em src/content.tex:152-156. APLICADA AO FONTE em 2026-08-21,
% por decisao do autor, para valer na versao final. O dissertacao.pdf entregue NAO foi
% reconstruido: ele continua sendo o registro do que a banca recebeu.
%
% O DEFEITO, em uma frase: o Resumo em portugues afirma que o modelo conjunto superou os
% dedicados na previsao da proxima categoria "em todos os conjuntos", e o resultado entregue
% e superioridade em UM conjunto apenas.
%
% POR QUE ISSO IMPORTA MAIS DO QUE O TAMANHO SUGERE. O Resumo e o primeiro paragrafo em
% portugues que um membro da banca le, e a frase generaliza exatamente na direcao que a
% WRITING_LAW §3 proibe: "outperforms" e reservado a superioridade pareada demonstrada, e
% uma diferenca que nao passa no teste e reportada pelo limite que seu intervalo sustenta,
% nunca como igualdade e nunca como resultado universal.
%
% O DEFEITO E ISOLADO -- medido, nao suposto. Todo o resto dos dois volumes esta correto e
% diz o oposto da frase do Resumo. Verificado por extracao de texto do dissertacao.pdf
% entregue (md5 5be69d1bf589e6fe5e794898a06306cf) em 2026-08-21:
%   Abstract (en), l. 748  : "it outperforms the dedicated model at Florida, and the five
%                             remaining differences are equivalent to zero within half a point"
%   Cap. 5, l. 3320        : mesma frase, mesma delimitacao
%   Cap. 6 / contribuicoes : "it outperforms at one dataset, and every one of the six
%                             differences is equivalent to zero within half a point"
%   §2.5.2                 : idem
% Ou seja: a versao em ingles do proprio Resumo ja esta certa. A correcao e alinhar o
% portugues ao ingles, nao mudar um resultado.
%
% A FONTE DO VEREDITO: src/tables/mobiwac/results.tex (celula FL categoria, 37.35 dedicado
% contra 37.55 conjunto, marcada com a seta de superioridade) e
% ../evidence/ladder_recompute.json (Holm p = 0.011 em Florida; as outras cinco nao resolvidas).
% ============================================================================

\paragraph*{Errata: o escopo do resultado de categoria no Resumo.}
No Resumo em portugues, o periodo que relata os resultados passa a ler:

\begin{quote}
Essas analises compararam o modelo conjunto com os modelos dedicados a cada tarefa: na
previsao da proxima regiao ele os superou nos dois conjuntos com os maiores vocabularios de
regiao e permaneceu dentro da margem de dois pontos de Acc@10, registrada antes de qualquer
resultado ser lido, nos outros quatro (nao-inferioridade estatistica, por dois testes
unilaterais); e na previsao da proxima categoria ele os superou em um conjunto, enquanto as
cinco diferencas restantes sao equivalentes a zero dentro de meio ponto.
\end{quote}

\noindent Este periodo e a traducao fiel do que o Abstract em ingles ja afirmava. A frase
anterior --- ``superou os modelos dedicados na previsao da proxima categoria em todos os
conjuntos'' --- generalizava um resultado que o proprio documento delimita a Florida, e nao
tinha equivalente em nenhuma outra passagem dos dois volumes.

% ----------------------------------------------------------------------------
% RESPOSTA ORAL, se um arguidor citar o Resumo entregue:
%
% "O senhor esta certo, e o erro e meu. Essa frase do Resumo em portugues generaliza um
%  resultado que o documento delimita: a superioridade em categoria e em Florida, com Holm
%  p = 0,011, e as outras cinco diferencas nao sao resolvidas, cada uma limitada a meio ponto
%  de zero. O Abstract em ingles, a Secao 2.5, o Capitulo 5 e o Capitulo 6 dizem exatamente
%  isso; o defeito ficou so na versao em portugues do Resumo, e ele entra como errata no
%  deposito final."
%
% Regra para a apresentacao: NUNCA dizer a frase do Resumo. O veredito falado e sempre
% "supera em Florida" para categoria. Se ninguem perguntar, nao levantar o assunto --- a
% errata ja cobre o deposito.
% ----------------------------------------------------------------------------
